\section*{Chapter \ref{ch:b2:project-key-value-server} --- Project: A Key-Value Server}

\begin{solution}{exo:b2:project-key-value-server:1}
\omoutput{code/bachelor-2/31-project-key-value-server/outputs/exo_session.txt}{The session's replies.}
The second \code{DEL b} finds nothing. \code{SET c 4 5} has three arguments and is refused. The line with doubled spaces
is accepted: the parser drops the \omterm{def:b2:finite-automata:word}{empty words} between spaces, so it reads as \code{SET c 4}. A stricter protocol would
refuse it; this one is lenient on purpose, and the tests fix the behaviour so that it cannot change by accident.
\code{DEL} without a key is refused, and the last \code{COUNT} sees \code{a} and \code{c}.
\end{solution}

\begin{solution}{exo:b2:project-key-value-server:2}
\omoutput{code/bachelor-2/31-project-key-value-server/outputs/exo_one_read.txt}{Three reads: two requests in the first,
half a request in each of the others.}
TCP delivers a stream of bytes, not the client's writes. The first \code{read} returns both requests of the first write,
which the buggy server would parse as one line, with a newline inside. The last request was written in two pieces, 100
milliseconds apart, and arrives in two reads, neither of them a request. \Cref{lst:b2:project-key-value-server:handle}
reads through a buffer up to each newline: it would see \code{SET a 1}, \code{GET a} and \code{GET a}, whatever the
reads return.
\end{solution}

\begin{solution}{exo:b2:project-key-value-server:3}
\omoutput{code/bachelor-2/31-project-key-value-server/outputs/exo_incr.txt}{Two increments make 1; eight clients of
INCR lose none.}
Both clients read 0 before either writes, then both write 1: a \omterm{def:b2:database-design:anomalies}{lost update}, staged here by a \omterm{def:b2:threads-and-synchronisation:barrier}{barrier} between the read
and the write. Each step is locked, but the \omterm{def:b2:threads-and-synchronisation:lock}{lock} is released between them. \code{INCR} reads, adds one and writes under
one acquisition of the \omterm{def:b2:threads-and-synchronisation:lock}{lock}, so no other client can come between: $8 \times 10\,000 = 80\,000$. In the server, it is one
more \code{Command} handled like \code{SET} under the shard's write \omterm{def:b2:threads-and-synchronisation:lock}{lock}, and logged as the \code{SET} of the new value,
so that a \omterm{def:b2:project-key-value-server:log}{replay} does not depend on the order of increments being repeated exactly.
\end{solution}

\begin{solution}{exo:b2:project-key-value-server:4}
A record is $8 + 3 + 1 + 4 + 1 + 6 = 23$ bytes. Before compaction the log holds $23n$ bytes; after it, one record per
key present, $23\min(n, k)$ bytes, here $23k$ since $n \ge k$.
\omoutput{code/bachelor-2/31-project-key-value-server/outputs/exo_log_size.txt}{The four cases.}
When $n = k$ compaction gains nothing; with $n = 100\,000$ and $k = 10$ it shrinks the log ten thousand times.
\end{solution}

\begin{solution}{exo:b2:project-key-value-server:5}
Each client's two steps keep their order, so an \omterm{def:b2:threads-and-synchronisation:interleaving}{interleaving} is a choice of the 2 positions of $a$'s steps among 4:
$\binom{4}{2} = 6$.
\omoutput{code/bachelor-2/31-project-key-value-server/outputs/exo_interleave.txt}{The six \omterm{def:b2:threads-and-synchronisation:interleaving}{interleavings}.}
Only the two serial orders give 2. In the four others both clients read 0 before either writes, and both write 1: one
increment is lost. That is the \omterm{def:b2:database-design:anomalies}{lost update} of the previous chapter, here with no database to refuse it.
\end{solution}

\begin{solution}{exo:b2:project-key-value-server:6}
\omoutput{code/bachelor-2/31-project-key-value-server/outputs/exo_mset_deadlock.txt}{Opposite orders deadlock; one
order does not.}
The wait-for graph has an edge from \code{MSET a b} to \code{MSET b a} (the first waits for shard 1, held by the second)
and an edge back (the second waits for shard 0, held by the first): a cycle, so a deadlock; the program sees it only
because it gives up after 200 milliseconds. Fix: sort the shard numbers of the keys and \omterm{def:b2:threads-and-synchronisation:lock}{lock} in increasing order, as in
\cref{prop:b2:project-key-value-server:lock-order}; both now take shard 0 first, so one of them waits before holding
anything, and the cycle cannot close. Two keys in the same shard need one \omterm{def:b2:threads-and-synchronisation:lock}{lock}, not two: the numbers must be deduplicated,
or a \omterm{def:b2:threads-and-synchronisation:thread}{thread} would wait for a \omterm{def:b2:threads-and-synchronisation:lock}{lock} it already holds.
\end{solution}

\begin{solution}{exo:b2:project-key-value-server:7}
\omoutput{code/bachelor-2/31-project-key-value-server/outputs/exo_setex.txt}{Lazy expiry on read, then a sweep.}
The key set at $t = 0$ for 30 seconds is still there at 29 and gone at 30. Two keys were never read after expiring:
only the sweep removes them, and a real server runs it from time to time. The log must record the instant of expiry, not
the duration: a \code{SETEX k 30 v} replayed an hour later would give the key thirty new seconds. Record
\code{SETAT k} \emph{deadline} \code{v}, and let \omterm{def:b2:project-key-value-server:log}{replay} skip, or delete, the keys whose deadline has passed.
\end{solution}

\begin{solution}{exo:b2:project-key-value-server:8}
Write $M_0$ for the empty map and $M_i$ for the map after applying $r_1, \dots, r_i$. \emph{Claim}: after $i$ iterations
of the loop of \cref{lst:b2:project-key-value-server:replay}, the map is $M_i$ and \code{at} is the offset where
$r_{i+1}$ begins. True for $i = 0$. If true for $i < n$, \code{next_record} reads $r_{i+1}$: its eight header bytes are
present, its length is right, its CRC matches, since the record was written whole by \code{record}, and its payload
parses as a \code{SET} or \code{DEL}; applying it gives $M_{i+1}$, and \code{at} advances by its size. At $i = n$, the
bytes that remain do not begin with a complete, undamaged record by hypothesis, so \code{next_record} returns
\code{None} and the loop ends with $M_n$. Finally, the server appended each change before applying it, in the order of
the changes to each key (the shard \omterm{def:b2:threads-and-synchronisation:lock}{lock} is held across both), and changes to different keys commute; so $M_n$ is the map
the server held after its $n$-th logged change. The hypothesis on the tail is the delicate one: a torn record is short
and fails the length test, and a tail of zeros fails the parse; random bytes pass the CRC test with probability $2^{-32}$.
\end{solution}

\begin{solution}{exo:b2:project-key-value-server:9}
\omoutput{code/bachelor-2/31-project-key-value-server/outputs/exo_snapshot.txt}{Snapshot and tail replayed as one log.}
The snapshot is itself a log of one \code{SET} per key, so recovery \omterm{def:b2:project-key-value-server:log}{replays} the snapshot and then the tail as one
sequence; a \code{DEL} in the tail removes a key of the snapshot, as \code{k4} shows. Order: write the snapshot to a
temporary file, sync it, rename it over the old snapshot, sync the directory, and only then empty the log. A crash while
the snapshot is written leaves the temporary file, which recovery ignores, and the old snapshot and full log, which are
still correct. Emptying the log before the snapshot is durable would lose every change since the previous snapshot. The
program does not sync, since it demonstrates the formats only; a server must.
\end{solution}

\begin{solution}{exo:b2:project-key-value-server:10}
\omoutput{code/bachelor-2/31-project-key-value-server/outputs/exo_mset.txt}{The two keys in different \omterm{def:b2:project-key-value-server:shard}{shards}, and no torn
read.}
The module \code{mset.rs} sorts and deduplicates the shard numbers, takes the write \omterm{def:b2:threads-and-synchronisation:lock}{locks} in that order, and changes
every key while it holds them all; \code{mget} takes the read \omterm{def:b2:threads-and-synchronisation:lock}{locks} in the same order. An \code{MGET} therefore sees the
state either before or after each \code{MSET}, never between, and no deadlock is possible
(\cref{prop:b2:project-key-value-server:lock-order}), although half the writers name the keys the other way round. The
crate's test runs the same check with three keys. In the server, an \code{MSET} must also be one log record, or a crash
between its records would \omterm{def:b2:project-key-value-server:log}{replay} half of it.
\end{solution}

\begin{solution}{exo:b2:project-key-value-server:11}
Each client sends a request every $s$ microseconds, so $c$ clients offer $c / s$ requests per microsecond. One \omterm{def:b2:threads-and-synchronisation:lock}{lock} is
held $h$ microseconds per request, so it lets through at most $1 / h$ requests per microsecond, whatever $c$; $N$ \omterm{def:b2:project-key-value-server:shard}{shards},
with keys spread evenly, at most about $N / h$. The throughput is about $\min(c / s,\ N / h)$. In process, one \omterm{def:b2:threads-and-synchronisation:thread}{thread} does
about 4.5 million operations per second, so $s \approx 0.2$~\textmu s, and the \omterm{def:b2:threads-and-synchronisation:lock}{lock} is held for most of it: one \omterm{def:b2:threads-and-synchronisation:lock}{lock} caps
near 5 million per second. The measured throughput \emph{drops} below that with more \omterm{def:b2:threads-and-synchronisation:thread}{threads}, to 1 million: the model
leaves out the cost of contention itself, the \omterm{def:b2:threads-and-synchronisation:lock}{lock}'s memory moving between cores and the \omterm{def:b2:threads-and-synchronisation:thread}{threads} put to sleep and woken.
Over TCP, $s \approx 150$~\textmu s, so 8 clients offer about 53 thousand requests per second, far below either cap,
and the \omterm{def:b2:threads-and-synchronisation:lock}{lock} does not matter, as measured. The measured rate at 8 clients is higher than the model, about 100 thousand,
because the latency per request fell as the machine got busier; a model with one fixed $s$ cannot show that.
\end{solution}

\begin{solution}{exo:b2:project-key-value-server:12}
\omoutput{code/bachelor-2/31-project-key-value-server/outputs/exo_pipeline.txt}{One client, batches of 1 to 1000
requests, on the author's laptop.}
With $b = 1$ each request pays a whole round trip, about 150 microseconds of \omterm{def:b2:processes-and-system-calls:syscall}{system calls}, scheduling and loopback. A
batch pays one round trip for $b$ requests, so the time per request falls almost as $1/b$ until the server's own work
per request, under a microsecond, dominates: about 230 times the throughput at $b = 1000$. \Cref{lst:b2:project-key-value-server:handle}
reads requests from its buffer one line at a time whatever the batch, and \omterm{def:b2:files-pipes-signals:buffering}{flushes} the replies only when the buffer is
empty: one write of $b$ replies for a batch of $b$ requests. Nothing in the protocol had to change, because each reply
is one line in the order of the requests.
\end{solution}
